\documentclass[11pt,a4paper]{article}

% ----------------------------------------------------------------------
% Safety Case Report
% Introduction to Machine Learning Safety (SoSe 2026)
% Otto-von-Guericke University Magdeburg
% ----------------------------------------------------------------------

\usepackage[utf8]{inputenc}
\usepackage[T1]{fontenc}
\usepackage[margin=2.5cm]{geometry}
\usepackage{amsmath,amssymb}
\usepackage{booktabs}
\usepackage{array}
\usepackage{longtable}
\usepackage{graphicx}
\usepackage[table]{xcolor}
\usepackage{enumitem}
\usepackage{titlesec}
\usepackage{fancyhdr}
\usepackage{tcolorbox}
\usepackage[hidelinks]{hyperref}

\definecolor{claimblue}{RGB}{30,75,140}
\definecolor{warnred}{RGB}{170,30,30}
\definecolor{lightgray}{RGB}{240,240,240}

\titleformat{\section}{\large\bfseries\color{claimblue}}{\thesection}{0.6em}{}
\titleformat{\subsection}{\normalsize\bfseries}{\thesubsection}{0.6em}{}

\pagestyle{fancy}
\fancyhf{}
\lhead{\small Safety Case Report}
\rhead{\small CARLA Perception System}
\cfoot{\thepage}
\renewcommand{\headrulewidth}{0.4pt}

% Claim / evidence boxes -------------------------------------------------
\newtcolorbox{claimbox}[1]{colback=lightgray,colframe=claimblue,
  fonttitle=\bfseries,title=#1,boxrule=0.6pt,arc=2pt,left=6pt,right=6pt,top=4pt,bottom=4pt}

\newcommand{\claim}[2]{\vspace{2pt}\noindent\textbf{\color{claimblue}#1:} #2\par\vspace{2pt}}

\begin{document}

% ======================================================================
% TITLE
% ======================================================================
\begin{titlepage}
\centering
\vspace*{1.5cm}
{\Huge\bfseries Safety Case Report\par}
\vspace{0.6cm}
{\LARGE CARLA Autonomous-Driving Perception System\par}
\vspace{0.4cm}
{\large A Structured Safety Argument for a ResNet-18 Based\\ Pedestrian, Vehicle and Traffic-Light Detection Stack\par}
\vspace{2.0cm}

\begin{tabular}{rl}
\textbf{Course:} & Introduction to Machine Learning Safety \\
\textbf{Institution:} & Otto-von-Guericke University Magdeburg (OvGU) \\
\textbf{Semester:} & Summer Term 2026 \\
\textbf{Author:} & Vivek Chitradurga Mallikarjun \\
\textbf{Document type:} & Assurance / Safety Case \\
\textbf{Status:} & Draft for review \\
\end{tabular}

\vfill
{\small This report presents a structured safety case for a machine-learning
perception system. All quantitative evidence is drawn from the experimental
exercises (training, evaluation, backdoor analysis, explainability,
out-of-distribution detection, adversarial robustness and uncertainty
calibration) performed on the CARLA dataset. The report is intended for
educational purposes only and does not certify the system for real-world
deployment.\par}
\vspace{1cm}
{\small \today}
\end{titlepage}

\tableofcontents
\newpage

% ======================================================================
\section{Executive Summary}
% ======================================================================
This document constructs a safety case for a camera-based perception system
composed of three independent binary classifiers --- \emph{pedestrian
detection}, \emph{vehicle detection} and \emph{traffic-light detection} ---
each built on an ImageNet-pretrained ResNet-18 backbone and trained on images
from the CARLA driving simulator.

A safety case is a \emph{structured, evidence-supported argument} that a system
is acceptably safe for a given use within a defined operating context. Following
that principle, this report states an explicit top-level safety goal, decomposes
it into verifiable sub-claims, and links each sub-claim to concrete evidence
produced by the course exercises.

\begin{tcolorbox}[colback=warnred!8,colframe=warnred,title=\bfseries Overall verdict]
The evidence gathered does \textbf{not} support an unconditional claim of
acceptable safety. The pedestrian classifier --- the most safety-critical
component --- exhibits a test-set recall of only \textbf{0.19}, all three
models are severely \textbf{over-confident}, and the system is vulnerable to
distribution shift, backdoor poisoning and adversarial perturbation. The system
must therefore be regarded as \textbf{not fit for deployment} in its current
form. Section~\ref{sec:recommendations} lists the conditions that would have to
be met before deployment could be reconsidered.
\end{tcolorbox}

% ======================================================================
\section{System Description and Operational Design Domain}
% ======================================================================
\subsection{System under assessment}
The item under assessment is the perception front-end of a CARLA autonomous
driving stack. It ingests front-facing RGB camera frames and produces three
independent binary decisions:

\begin{center}
\begin{tabular}{lll}
\toprule
\textbf{Model} & \textbf{Output} & \textbf{Backbone} \\
\midrule
Pedestrian detector    & \texttt{has\_pedestrian} $\in \{0,1\}$    & ResNet-18 (ImageNet-pretrained) \\
Vehicle detector       & \texttt{has\_vehicle} $\in \{0,1\}$       & ResNet-18 (ImageNet-pretrained) \\
Traffic-light detector & \texttt{has\_traffic\_light} $\in \{0,1\}$& ResNet-18 (ImageNet-pretrained) \\
\bottomrule
\end{tabular}
\end{center}

Each model applies a sigmoid to a single logit and thresholds the resulting
probability (default $\tau = 0.5$) to obtain a class decision. Inputs are
resized to $224\times224$ and normalised with ImageNet statistics.

\subsection{Operational Design Domain (ODD)}
The ODD, derived from the training distribution and refined in the
out-of-distribution study, is:

\begin{itemize}[leftmargin=1.4em,itemsep=1pt]
  \item \textbf{Environment:} daytime, clear-weather urban driving.
  \item \textbf{Scene:} standard traffic, moderate vehicle speed, urban road layouts.
  \item \textbf{Sensor:} clean, unobstructed front-facing RGB camera.
  \item \textbf{Excluded conditions:} fog, night, heavy precipitation, and
        previously unseen towns/geographies. These constitute
        out-of-distribution (OOD) inputs and are \emph{outside} the validated ODD.
\end{itemize}

A key finding of the OOD study is that the original ODD specification was
incomplete: it constrained weather, lighting and road type but did not state
whether \emph{unseen towns} are admissible. This report treats geographic
generalisation as \emph{out of ODD} until validated.

% ======================================================================
\section{Hazard Analysis}
% ======================================================================
Because the perception output feeds downstream planning and control, a
perception error can propagate into an unsafe vehicle action. The dominant
hazards and their asymmetry are:

\begin{longtable}{p{2.6cm} p{4.6cm} p{2.2cm} p{3.4cm}}
\toprule
\textbf{Hazard} & \textbf{Perception failure} & \textbf{Severity} & \textbf{Dominant error mode} \\
\midrule
\endhead
H1 --- Pedestrian collision & Missed pedestrian (false negative) & Catastrophic & Recall failure \\
H2 --- Vehicle collision & Missed vehicle (false negative) & Critical & Recall failure \\
H3 --- Traffic-violation / junction collision & Missed / misread traffic light & Critical & Recall / class error \\
H4 --- Unnecessary emergency braking & Phantom detection (false positive) & Marginal & Precision failure \\
H5 --- Silent failure on OOD input & Confident but wrong output off-ODD & Critical & Confidence \& shift \\
H6 --- Malicious manipulation & Backdoor / adversarial evasion & Catastrophic & Security \\
\bottomrule
\end{longtable}

\begin{claimbox}{Key hazard-analysis conclusion}
For H1--H3 a \textbf{false negative} (missing a real object) is far more
dangerous than a false positive. The safety argument must therefore be driven
primarily by \textbf{recall} and by the \textbf{calibrated confidence} of the
detectors, not by aggregate accuracy, which is misleading under class imbalance.
\end{claimbox}

% ======================================================================
\section{Safety Argument Structure}
% ======================================================================
The safety case is organised as a goal-decomposition argument (in the spirit of
Goal Structuring Notation).

\begin{claimbox}{G0 --- Top-level safety goal}
The CARLA perception system provides object detections that are sufficiently
reliable, robust, interpretable and calibrated to be used safely within its
declared ODD.
\end{claimbox}

\noindent G0 is decomposed into six sub-goals, each supported by one exercise:

\begin{itemize}[leftmargin=1.6em,itemsep=2pt]
  \item \textbf{G1 (Nominal performance):} detectors meet minimum per-class
        performance thresholds on in-ODD data. \hfill\emph{Evidence: Ex.~4}
  \item \textbf{G2 (Data-supply-chain integrity):} the training pipeline is not
        compromised by data poisoning. \hfill\emph{Evidence: Ex.~5}
  \item \textbf{G3 (Interpretability):} predictions are driven by causally
        relevant image regions. \hfill\emph{Evidence: Ex.~6}
  \item \textbf{G4 (OOD awareness):} the system detects inputs outside its ODD.
        \hfill\emph{Evidence: Ex.~7}
  \item \textbf{G5 (Adversarial robustness):} performance degrades gracefully
        under bounded input perturbation. \hfill\emph{Evidence: Ex.~8}
  \item \textbf{G6 (Trustworthy uncertainty):} reported confidence is calibrated
        and supports risk-aware decisions. \hfill\emph{Evidence: Ex.~9}
\end{itemize}

The following sections examine the evidence for each sub-goal and record whether
it is \textbf{satisfied}, \textbf{partially satisfied} or \textbf{refuted}.

% ======================================================================
\section{Evidence G1 --- Nominal Performance (Exercise 4)}
% ======================================================================
Per-class performance was measured on the in-ODD test split using precision,
recall and F1-score.

\begin{center}
\begin{tabular}{lccc}
\toprule
\textbf{Model} & \textbf{Precision} & \textbf{Recall} & \textbf{F1-score} \\
\midrule
Vehicle        & 0.986 & 0.800 & 0.883 \\
\rowcolor{warnred!12}
Pedestrian     & 0.572 & \textbf{0.191} & 0.287 \\
Traffic light  & 0.916 & 0.991 & 0.952 \\
\bottomrule
\end{tabular}
\end{center}

\claim{Finding}{The \textbf{pedestrian} model has by far the lowest recall
($0.191$): it misses roughly four out of every five pedestrians actually present.
Because H1 (pedestrian collision) is the most severe hazard, this is a
first-order safety defect. The vehicle model's recall ($0.80$) is moderate; the
traffic-light model performs well ($\mathrm{F1}=0.95$).}

\claim{Safety threshold argument}{For a safety-critical pedestrian detector, a
recall floor of at least $\mathbf{0.95}$ (ideally $\ge 0.99$) is a reasonable
pre-deployment requirement, since every missed pedestrian corresponds directly
to a potential catastrophic outcome and false positives (unnecessary braking)
are comparatively benign. The measured $0.191$ is nearly five times below this
floor.}

\begin{tcolorbox}[colback=warnred!8,colframe=warnred]
\textbf{Verdict on G1: REFUTED.} The pedestrian detector does not meet any
defensible minimum recall for a safety-critical function.
\end{tcolorbox}

% ======================================================================
\section{Evidence G2 --- Data-Supply-Chain Integrity (Exercise 5)}
% ======================================================================
A backdoor (data-poisoning) attack was implemented against the pedestrian model.
A small red $10\times10$ pixel patch in the top-left corner served as the
trigger. During training, $10\%$ of positive (pedestrian-present) images were
stamped with the trigger and had their label flipped to ``no pedestrian''.

\begin{itemize}[leftmargin=1.4em,itemsep=1pt]
  \item \textbf{Attack goal:} force the model to report \emph{no pedestrian}
        whenever the trigger is present, while behaving normally otherwise.
  \item \textbf{Metrics:} clean recall (utility on untriggered data) and
        \emph{attack success rate} (fraction of triggered pedestrian images
        classified as ``no pedestrian'').
\end{itemize}

\claim{Finding}{The experiment demonstrates that a low poisoning rate ($10\%$ of
one class) is sufficient to install a stealthy backdoor: the poisoned model can
retain acceptable clean-data behaviour while failing on any triggered input.
This directly realises hazard H6 and shows that the training data supply chain is
a genuine attack surface.}

\begin{tcolorbox}[colback=warnred!8,colframe=warnred]
\textbf{Verdict on G2: NOT SATISFIED.} No data-provenance, poisoning-detection or
trigger-scanning control is present. Integrity of the training set cannot be
claimed.
\end{tcolorbox}

% ======================================================================
\section{Evidence G3 --- Interpretability (Exercise 6)}
% ======================================================================
Grad-CAM was used to visualise the image regions most responsible for each
prediction, for both correct predictions and misclassifications, and across
nominal and shifted conditions (clean, fog, night, alternative town).

\claim{Purpose}{Interpretability is used here as \emph{supporting} rather than
\emph{primary} evidence: it helps diagnose \emph{why} the quantitative failures
occur --- e.g.\ whether the pedestrian model attends to spurious background
regions rather than to pedestrians, and how attention maps degrade under
distribution shift.}

\claim{Finding}{Grad-CAM heatmaps provide qualitative insight into failure modes
and are a valuable debugging and transparency tool. They are, however,
explanatory rather than a guarantee of correctness: a plausible heatmap does not
by itself certify a safe decision.}

\begin{tcolorbox}[colback=orange!8,colframe=orange!70!black]
\textbf{Verdict on G3: PARTIALLY SATISFIED.} Interpretability tooling exists and
supports diagnosis, but does not on its own discharge any safety claim.
\end{tcolorbox}

% ======================================================================
\section{Evidence G4 --- Out-of-Distribution Awareness (Exercise 7)}
% ======================================================================
Robustness to distribution shift was studied across three OOD domains --- fog and
night (appearance shift) and an unseen town (structural/domain shift) --- and two
OOD-detection methods were compared.

\subsection{Behaviour under shift}
Reported per-model accuracy under each condition:

\begin{center}
\begin{tabular}{lcccc}
\toprule
\textbf{Model} & \textbf{In-ODD} & \textbf{Fog} & \textbf{Night} & \textbf{Unseen town} \\
\midrule
Pedestrian     & 0.945 & 0.976 & 0.976 & 0.958 \\
Vehicle        & 0.904 & \textbf{0.609} & 0.862 & --- \\
Traffic light  & 0.963 & 0.814 & 1.000 & --- \\
\bottomrule
\end{tabular}
\end{center}

\claim{Caveat}{These accuracy figures are inflated by class imbalance (see
Exercise~9) and must not be read as evidence of safety. The vehicle model
nevertheless shows a large, unambiguous accuracy drop under fog
($0.90 \rightarrow 0.61$), confirming appearance shift as a real hazard (H5).}

\subsection{OOD detector comparison}
\begin{center}
\begin{tabular}{lc}
\toprule
\textbf{OOD detection method} & \textbf{AUROC} \\
\midrule
Maximum Softmax Probability (MSP) & 0.751 \\
\rowcolor{claimblue!12}
Feature-space $k$-NN              & \textbf{0.983} \\
\bottomrule
\end{tabular}
\end{center}

\claim{Finding}{Confidence-based MSP is a weak OOD detector (AUROC $0.75$),
consistent with the models' over-confidence. Feature-space $k$-NN, which measures
distance to training embeddings, is far more reliable (AUROC $0.98$) and is the
recommended runtime OOD monitor.}

\begin{tcolorbox}[colback=orange!8,colframe=orange!70!black]
\textbf{Verdict on G4: PARTIALLY SATISFIED (with condition).} A strong OOD
detector ($k$-NN) is available and \emph{must} be deployed as an input-validity
gate. MSP alone is inadequate.
\end{tcolorbox}

% ======================================================================
\section{Evidence G5 --- Adversarial Robustness (Exercise 8)}
% ======================================================================
Robustness was probed with the Fast Gradient Sign Method (FGSM), which perturbs
an input in the direction that maximises the loss:
\[
  x_{\mathrm{adv}} = x + \epsilon \cdot \operatorname{sign}\!\big(\nabla_x J(\theta, x, y)\big).
\]
Increasing the perturbation budget $\epsilon$ measures how quickly accuracy and
confidence collapse.

\claim{Finding}{As with essentially all undefended CNN classifiers, small,
visually imperceptible perturbations are sufficient to flip predictions and
sharply degrade accuracy. No adversarial training or input pre-processing defence
is present, so hazard H6 (malicious manipulation) is unmitigated. Combined with
the backdoor result (G2), the system has no security assurance against an active
adversary.}

\begin{tcolorbox}[colback=warnred!8,colframe=warnred]
\textbf{Verdict on G5: NOT SATISFIED.} The models are undefended against bounded
adversarial perturbations.
\end{tcolorbox}

% ======================================================================
\section{Evidence G6 --- Trustworthy Uncertainty (Exercise 9)}
% ======================================================================
Calibration quality was measured with the Expected Calibration Error (ECE),
\[
  \mathrm{ECE} = \sum_{m=1}^{M}\frac{|B_m|}{n}\,\big|\operatorname{acc}(B_m)-\operatorname{conf}(B_m)\big|,
\]
followed by temperature scaling and a cost-sensitive decision-threshold analysis.

\subsection{Calibration before/after temperature scaling}
\begin{center}
\begin{tabular}{lccc}
\toprule
\textbf{Model} & \textbf{ECE (raw)} & \textbf{Best $T$} & \textbf{ECE (scaled)} \\
\midrule
Pedestrian     & 0.156 & 2.9 & 0.056 \\
Vehicle        & 0.736 & 3.0 & 0.557 \\
Traffic light  & 0.597 & 3.0 & 0.394 \\
\bottomrule
\end{tabular}
\end{center}

\claim{Finding}{All three models are strongly \textbf{over-confident} (mean
confidence far exceeds mean accuracy; e.g.\ the vehicle model reports mean
confidence $0.99$ against mean accuracy $0.25$). Temperature scaling
($T>1$ for every model) reduces ECE substantially for the pedestrian model
($0.156 \rightarrow 0.056$) but only partially for the vehicle and traffic-light
models, whose residual error reflects genuine accuracy deficits that calibration
alone cannot fix.}

\subsection{Cost-sensitive decision making}
Using an asymmetric cost of \textbf{1} for a false positive and \textbf{100} for
a false negative (pedestrian model), the total decision loss was compared at the
default threshold $\tau=0.5$ and at a cost-optimal threshold $\tau^{\ast}$:

\begin{center}
\begin{tabular}{lccc}
\toprule
\textbf{Configuration} & \textbf{Threshold} & \textbf{FP / FN} & \textbf{Total cost} \\
\midrule
Default decision     & $\tau=0.5$        & $0 \,/\, 706$ & 70{,}600 \\
Uncalibrated + cost-opt & $\tau^{\ast}$   & $2765 \,/\, 5$ & 3{,}265 \\
\rowcolor{claimblue!12}
Calibrated + cost-opt   & $\tau^{\ast}$   & $2894 \,/\, 0$ & \textbf{2{,}894} \\
\bottomrule
\end{tabular}
\end{center}

\claim{Finding}{At the naive $\tau=0.5$ the pedestrian model predicts
\emph{negative for every input}, missing all $706$ pedestrians (recall $0$) --- a
catastrophic silent failure. Combining temperature scaling with a
cost-optimal threshold drives false negatives to \textbf{zero} (recall $1.0$) and
reduces the safety-weighted loss from $70{,}600$ to $2{,}894$, at the price of
many false positives (unnecessary braking events).}

\begin{tcolorbox}[colback=orange!8,colframe=orange!70!black]
\textbf{Verdict on G6: PARTIALLY SATISFIED (with condition).} Raw confidence is
untrustworthy, but calibration plus risk-aware thresholding restores usable, safe
decision behaviour \emph{for the pedestrian model}. This mitigation is mandatory,
not optional.
\end{tcolorbox}

% ======================================================================
\section{Consolidated Evidence Summary}
% ======================================================================
\begin{center}
\begin{tabular}{clc}
\toprule
\textbf{Goal} & \textbf{Claim} & \textbf{Verdict} \\
\midrule
G1 & Nominal per-class performance      & \textcolor{warnred}{\textbf{Refuted}} \\
G2 & Training-data integrity            & \textcolor{warnred}{\textbf{Not satisfied}} \\
G3 & Interpretability                   & \textcolor{orange!80!black}{\textbf{Partial}} \\
G4 & OOD awareness                      & \textcolor{orange!80!black}{\textbf{Partial (conditional)}} \\
G5 & Adversarial robustness             & \textcolor{warnred}{\textbf{Not satisfied}} \\
G6 & Trustworthy uncertainty            & \textcolor{orange!80!black}{\textbf{Partial (conditional)}} \\
\midrule
\multicolumn{2}{r}{\textbf{Top-level goal G0}} & \textcolor{warnred}{\textbf{Not met}} \\
\bottomrule
\end{tabular}
\end{center}

\claim{Argument conclusion}{No sub-goal is fully and unconditionally satisfied,
and two of the most safety-relevant sub-goals (G1 nominal performance and G5
adversarial robustness) are refuted or unmitigated. The top-level safety goal G0
therefore \textbf{cannot be discharged}.}

% ======================================================================
\section{Recommendations and Conditions for Deployment}
\label{sec:recommendations}
% ======================================================================
The following conditions are considered necessary (not necessarily sufficient)
before any deployment could be reconsidered:

\begin{enumerate}[leftmargin=1.6em,itemsep=2pt]
  \item \textbf{Raise pedestrian recall (blocks G1).} Retrain with rebalanced /
        augmented data, hard-negative and hard-positive mining, and a
        recall-oriented objective until test recall $\ge 0.95$.
  \item \textbf{Mandate calibrated, cost-aware decisions (G6).} Deploy
        temperature scaling and a cost-optimal threshold ($\tau^{\ast}$) for every
        safety-critical detector; never use the naive $\tau=0.5$.
  \item \textbf{Deploy a runtime OOD gate (G4).} Use the feature-space $k$-NN
        monitor (AUROC $0.98$) to reject or hand-over on off-ODD inputs (fog,
        night, unseen towns); do not rely on MSP.
  \item \textbf{Secure the data supply chain (G2).} Add data-provenance controls,
        poisoning/backdoor scanning and trigger-detection to the training pipeline.
  \item \textbf{Add adversarial defences (G5).} Introduce adversarial training
        and/or input sanitisation, and evaluate under FGSM and stronger attacks.
  \item \textbf{Complete and formalise the ODD.} Explicitly state admissible
        geographies, weather and lighting, and validate against them.
  \item \textbf{Adopt safety-relevant metrics.} Report recall and calibrated
        confidence per class; do not use aggregate accuracy as the headline
        metric under class imbalance.
\end{enumerate}

% ======================================================================
\section{Conclusion}
% ======================================================================
This safety case assembled evidence across nominal performance, data integrity,
interpretability, out-of-distribution awareness, adversarial robustness and
uncertainty calibration. The evidence shows that, while individual mitigations
(feature-space OOD detection, temperature scaling with cost-optimal thresholds)
are effective and should be mandatory, the system in its current form fails its
top-level safety goal --- principally because the pedestrian detector's recall is
far too low and because the system has no defence against poisoning or
adversarial manipulation. The honest conclusion of this assurance argument is a
\textbf{negative safety verdict}: the CARLA perception system is
\emph{not acceptably safe for deployment} until the conditions in
Section~\ref{sec:recommendations} are met and re-evidenced.

\vspace{0.8cm}
\noindent\rule{\linewidth}{0.4pt}\\
{\small \textit{Disclaimer.} This report was produced as coursework for the
\emph{Introduction to Machine Learning Safety} course at Otto-von-Guericke
University Magdeburg and is intended for educational and research purposes only.
It does not constitute a formal certification of the system.}

\end{document}
